\chapter{Modèles UML}

Les sources PlantUML et leurs exports SVG sont dans \code{docs/uml/}. Le cahier des charges (§8.2) prévoyait six
diagrammes ; trois diagrammes ont été ajoutés car ils décrivent des règles de gestion délicates (panier,
changement de statut) ou parce que la conception MVC était trop dense pour un seul diagramme.

\begin{table}[H]
  \centering\small
  \caption{Diagrammes UML du projet}
  \begin{tabularx}{\linewidth}{LP{4.6cm}P{2.6cm}}
    \ligneentete
    \entete{Diagramme} & \entete{Fichier} & \entete{Type} \\
    Cas d'utilisation Visiteur et Client & \code{cas-utilisation-client} & Cas d'utilisation \\
    Cas d'utilisation Administrateur & \code{cas-utilisation-admin} & Cas d'utilisation \\
    Classes du domaine & \code{classes-domaine} & Classes \\
    Conception MVC (1/2) : filtres, contrôleurs, services & \code{classes-conception-web} & Classes \\
    Conception MVC (2/2) : services, DAO, modèle & \code{classes-conception-metier} & Classes \\
    S'authentifier & \code{sequence-authentification} & Séquence \\
    Passer une commande & \code{sequence-passer-commande} & Séquence \\
    Ajouter au panier (ajout) & \code{sequence-ajouter-panier} & Séquence \\
    Changer le statut d'une commande (ajout) & \code{sequence-changer-statut} & Séquence \\
    \bottomrule
  \end{tabularx}
\end{table}

\section{Cas d'utilisation : Visiteur et Client}

Le Client hérite du Visiteur (généralisation). Le Visiteur consulte le catalogue (BF05), la fiche produit
(BF08), s'inscrit (BF01) et se connecte (BF02) ; le filtre par catégorie (BF06) et la recherche (BF07)
étendent la consultation du catalogue. Les cas du Client (panier, commande, historique) incluent
« Se connecter » (RG08). « Passer une commande » inclut « Consulter le panier » ; le détail d'une commande
(BF15) étend l'historique et l'annulation (BF16) étend le détail (figure~\ref{fig:cu-client}).

\diagramme{cas-utilisation-client}{Cas d'utilisation Visiteur et Client}{fig:cu-client}

\section{Cas d'utilisation : Administrateur}

Tous les cas du back-office incluent « Se connecter ». « Gérer » signifie créer, consulter, modifier et
supprimer. Le filtre par statut et le changement de statut étendent la consultation des commandes ;
l'activation ou la désactivation d'un compte étend la liste des utilisateurs. La légende rappelle les règles
RG04, RG13, RG14 et RG15 (figure~\ref{fig:cu-admin}).

\diagramme{cas-utilisation-admin}{Cas d'utilisation Administrateur}{fig:cu-admin}

\section{Classes du domaine}

Cinq entités persistantes : \code{Utilisateur}, \code{Categorie}, \code{Produit}, \code{Commande}, \code{LigneCommande}, et deux
énumérations : \code{Role} (CLIENT, ADMIN) et \code{StatutCommande} (EN\_ATTENTE, VALIDEE, EXPEDIEE, LIVREE, ANNULEE).
\code{LigneCommande} est une classe d'association porteuse d'attributs (quantité, prix unitaire) entre
\code{Commande} et \code{Produit}, à la place d'une association plusieurs-à-plusieurs simple. La relation
\code{Commande} – \code{LigneCommande} est une composition (1..* lignes). Les notes rappellent RG05, RG10, RG11, RG13
et RG15 (figure~\ref{fig:domaine}).

\diagramme{classes-domaine}{Diagramme de classes du domaine}{fig:domaine}

\begin{landscape}
\section{Conception MVC (1/2) : filtres, contrôleurs et services}

Ce diagramme (figure~\ref{fig:conception-web}) montre les trois filtres, leurs motifs d'URL et leur ordre
(déclarés dans \code{web.xml}, voir section~\ref{sec:securite}), les treize servlets avec leur URL
(\code{@WebServlet}), le service appelé par chacune et le renvoi vers les JSP de \code{WEB-INF/views}.
\code{TableauDeBordServlet} est la seule servlet qui appelle deux services (\code{CommandeService} et
\code{CatalogueService}), car le tableau de bord agrège des données des deux domaines.

\diagrammesuite{classes-conception-web}{Classes de conception MVC (1/2) : filtres, contrôleurs et services}{fig:conception-web}
\end{landscape}

\begin{landscape}
\section{Conception MVC (2/2) : services, DAO et modèle}

Signatures des méthodes des quatre services, de \code{RegleMetierException}, de \code{GenericDao<T>} et des quatre DAO
spécifiques, des classes du modèle (dont \code{Panier} et \code{LignePanier}, objets de session non persistés) et de
\code{HibernateUtil}, \code{DemarrageListener} et \code{Parametres}. Les noms de méthodes du code sont ceux de ce
diagramme (figure~\ref{fig:conception-metier}).

\diagrammesuite{classes-conception-metier}{Classes de conception MVC (2/2) : services, DAO et modèle}{fig:conception-metier}
\end{landscape}

\begin{landscape}
\section{Séquence : S'authentifier}

\code{ConnexionServlet} délègue à \code{AuthService.authentifier}, qui recherche l'utilisateur par e-mail
(HQL à paramètre nommé) et vérifie le mot de passe avec \code{BCrypt.checkpw}. Trois issues : identifiants
incorrects (même message que l'e-mail existe ou non), compte désactivé, ou succès avec redirection selon le
rôle (tableau de bord pour l'administrateur, catalogue pour le client). Au succès, l'ancienne session est
invalidée et une nouvelle session est créée (contre la fixation de session). Le diagramme
(figure~\ref{fig:seq-auth}) montre aussi la déconnexion, qui n'accepte que la méthode POST (un GET reçoit
l'erreur 405).

\diagrammesuite{sequence-authentification}{Diagramme de séquence : S'authentifier}{fig:seq-auth}
\end{landscape}

\begin{landscape}
\section{Séquence : Passer une commande}

Séquence centrale du projet (figure~\ref{fig:seq-commande}). Après le contrôle d'accès (\code{AuthFilter}),
\code{CommandeService.passerCommande} ouvre une transaction, recharge chaque produit du panier avec un verrou en
écriture (par identifiant croissant), vérifie qu'il est actif et que le stock suffit, décrémente le stock, crée
les lignes avec le prix courant (RG10), calcule le total (RG11) puis enregistre la commande et ses lignes par
cascade. Après le rechargement de l'utilisateur, un compte désactivé depuis sa connexion est refusé (BF22).
Toute anomalie provoque un rollback (RG12) et la page de validation est réaffichée avec le message, le panier
étant conservé.

\diagrammesuite{sequence-passer-commande}{Diagramme de séquence : Passer une commande}{fig:seq-commande}
\end{landscape}

\begin{landscape}
\section{Séquence : Ajouter au panier}

Le diagramme (figure~\ref{fig:seq-panier}) montre le refus par \code{AuthFilter} d'un visiteur ou d'un
administrateur (RG08), puis les contrôles de \code{PanierService} : produit inexistant ou inactif
(\code{ProduitDao.trouverActifParId}, RG15), quantité invalide, rupture (RG06), quantité cumulée supérieure au
stock (RG07). Le panier est un objet de la session HTTP ; le stock est revérifié à la validation de la commande.

\diagrammesuite{sequence-ajouter-panier}{Diagramme de séquence : Ajouter au panier}{fig:seq-panier}
\end{landscape}

\begin{landscape}
\section{Séquence : Changer le statut d'une commande}

Le diagramme (figure~\ref{fig:seq-statut}) montre d'abord le contrôle d'accès de \code{AdminFilter} :
redirection vers la connexion pour un visiteur, erreur 403 pour un client. \code{CommandeService.changerStatut}
verrouille la commande, vérifie que la transition est autorisée (RG13) et, pour une annulation, réintègre les
quantités dans le stock (RG14) par une requête \code{update} calculée par la base, dans la même transaction que
le changement de statut.

\diagrammesuite{sequence-changer-statut}{Diagramme de séquence : Changer le statut d'une commande}{fig:seq-statut}
\end{landscape}
